\documentclass[11pt,a4paper]{article}
\usepackage[margin=25mm]{geometry}
\usepackage{amsmath,amssymb,booktabs,array,siunitx,hyperref,xcolor,listings,enumitem,microtype}
\hypersetup{colorlinks=true,linkcolor=blue,urlcolor=blue,citecolor=blue,pdftitle={Report 1 -- Reservoir Model for OpenHPL/OpenHPLjl},pdfauthor={Madhusudhan Pandey}}
\lstdefinelanguage{Julia}{keywords={using,begin,end,function,return,if,else,elseif,for,while,let,struct,module,where,do,try,catch,finally,macro,local,global,const,export,import},sensitive=true,comment=[l]\#,morestring=[b]"}
\lstset{language=Julia,basicstyle=\ttfamily\small,breaklines=true,frame=single,columns=fullflexible,showstringspaces=false}
\title{Report 1 --- Reservoir Model for OpenHPL/OpenHPLjl}
\author{Madhusudhan Pandey}
\date{\today}
\begin{document}
\maketitle
\begin{abstract}
This report defines the reservoir as the first physical component in an equation-based hydropower modeling workflow using ModelingToolkit.jl and the SciML ecosystem. It follows the OpenHPLjl documentation sequence: updated context, generalized concept, literature model classes, OpenHPL mathematics and interpretation, Julia implementation, assembly, validation, expected outputs, and engineering interpretation.
\end{abstract}
\tableofcontents
\newpage

\section{Updated Context}
A reservoir is normally the upstream hydraulic boundary of a hydropower plant. It stores gravitational potential energy and establishes the hydraulic head available to the downstream waterway.

\[
\boxed{\text{Reservoir}\rightarrow\text{Intake}\rightarrow\text{Waterway}\rightarrow\text{Surge Tank}\rightarrow\text{Penstock}\rightarrow\text{Turbine}\rightarrow\text{Shaft}\rightarrow\text{Generator}\rightarrow\text{Grid}}
\]

OpenHPL is an open-source Modelica hydropower library containing waterway, electromechanical, controller, interface, and example models. The current OpenModelica-generated documentation identifies version 3.0.1 \cite{OpenHPLDocs}. Its user guide states that reservoirs, conduits, surge tanks, and fittings are based on mass and momentum balances, and that a hydrology model is available for reservoir inflow studies \cite{OpenHPLUserGuide}.

For short dynamic studies such as governor response or frequency containment reserve (FCR), the upper reservoir can often be approximated as a constant-head boundary. For longer simulations, reservoir storage dynamics should be retained.

The objective of the OpenHPLjl reservoir component is to provide a reusable equation-based model that can range from a constant-head reservoir to a nonlinear finite-storage reservoir without changing the downstream hydraulic interface.

\section{Generalized Concept Model}
The most general lumped reservoir representation is
\[
\boxed{\text{Inflow}\rightarrow\text{Storage}\rightarrow\text{Outflow}},
\qquad
Q_{\mathrm{in}}\rightarrow V(H)\rightarrow Q_{\mathrm{out}}.
\]

\begin{table}[ht]
\centering
\caption{Principal reservoir variables.}
\begin{tabular}{lll}
\toprule
Quantity & Symbol & SI unit \\
\midrule
Reservoir volume & $V$ & \si{\meter\cubed} \\
Water level / hydraulic head & $H$ & \si{\meter} \\
Inflow & $Q_{\mathrm{in}}$ & \si{\meter\cubed\per\second} \\
Outflow & $Q_{\mathrm{out}}$ & \si{\meter\cubed\per\second} \\
Free-surface area & $A$ & \si{\meter\squared} \\
Water density & $\rho$ & \si{\kilogram\per\meter\cubed} \\
Gravitational acceleration & $g$ & \si{\meter\per\second\squared} \\
\bottomrule
\end{tabular}
\end{table}

A generalized hydraulic connector contains
\[
\boxed{H,\quad Q},
\]
where $H$ is an across/potential variable and $Q$ is a through/flow variable. For two connected components,
\[
H_1=H_2,\qquad Q_1+Q_2=0.
\]

\section{Other Models Available in the Literature}
\subsection{Infinite or Constant-Head Reservoir}
\[
H(t)=H_0.
\]
There is no dynamic reservoir state. This is useful for short studies dominated by turbine, governor, waterway, or grid dynamics.

\subsection{Lumped Constant-Area Reservoir}
For a finite reservoir with approximately constant area,
\[
V=A_rH,
\qquad
A_r\frac{dH}{dt}=Q_{\mathrm{in}}-Q_{\mathrm{out}}.
\]

\subsection{Nonlinear Geometry Reservoir}
For realistic geometry,
\[
V=V(H),\qquad
A(H)=\frac{dV}{dH},
\]
so
\[
A(H)\dot H=Q_{\mathrm{in}}-Q_{\mathrm{out}}.
\]
The function $V(H)$ can come from surveyed geometry, fitted polynomials, or lookup tables.

\subsection{Reservoir With Hydrological Dynamics}
A longer-term model can use
\[
Q_{\mathrm{in}}=Q_{\mathrm{rain}}+Q_{\mathrm{snow}}+Q_{\mathrm{catchment}}-Q_{\mathrm{evaporation}}.
\]
OpenHPL includes hydrological modeling capability for reservoir inflow \cite{OpenHPLUserGuide}.

\subsection{Distributed or Spatial Reservoir Models}
For large lakes or flood studies, water level may vary spatially and PDE-based shallow-water or free-surface equations may be required.

A useful hierarchy is
\[
\boxed{\text{Constant head}\rightarrow\text{Lumped storage}\rightarrow\text{Nonlinear geometry}\rightarrow\text{Hydrological reservoir}\rightarrow\text{Distributed PDE model}}.
\]

\section{Mathematics Involved in the OpenHPL Concept/Model}
\subsection{Fundamental Conservation Law}
For incompressible water,
\[
\frac{dm}{dt}=\dot m_{\mathrm{in}}-\dot m_{\mathrm{out}},\qquad m=\rho V.
\]
With approximately constant density,
\[
\boxed{\frac{dV}{dt}=Q_{\mathrm{in}}-Q_{\mathrm{out}}}.
\]

\subsection{Constant-Area Reservoir}
If $V=A_rH$,
\[
\boxed{A_r\dot H=Q_{\mathrm{in}}-Q_{\mathrm{out}}},
\qquad
\boxed{\dot H=\frac{Q_{\mathrm{in}}-Q_{\mathrm{out}}}{A_r}}.
\]
The dynamic state is $x=H$.

\subsection{Nonlinear Reservoir Geometry}
For $V=V(H)$,
\[
\dot V=\frac{dV}{dH}\dot H.
\]
With $A(H)=dV/dH$,
\[
\boxed{A(H)\dot H=Q_{\mathrm{in}}-Q_{\mathrm{out}}}.
\]
For example,
\[
A(H)=a_0+a_1H+a_2H^2
\]
gives
\[
\dot H=\frac{Q_{\mathrm{in}}-Q_{\mathrm{out}}}{a_0+a_1H+a_2H^2}.
\]

\subsection{Hydraulic Pressure and Head}
At an outlet elevation $z_o$,
\[
p_o=p_{\mathrm{atm}}+\rho g(H-z_o).
\]
Hydraulic head can be written
\[
H_p=\frac{p}{\rho g}+z.
\]

\subsection{State-Space / DAE Form}
A convenient organization is
\[
x=H,
\qquad
z=\begin{bmatrix}V\\A\end{bmatrix},
\qquad
u=\begin{bmatrix}Q_{\mathrm{in}}\\Q_{\mathrm{out}}\end{bmatrix}.
\]
The equation-based form is
\[
\dot x=f(x,z,u,p),\qquad 0=g(x,z,u,p),
\]
with
\[
\dot H=\frac{Q_{\mathrm{in}}-Q_{\mathrm{out}}}{A(H)},\qquad V=V(H),\qquad A=\frac{dV}{dH}.
\]

\section{OpenHPL-Specific Modeling Interpretation}
OpenHPL is a component-oriented Modelica library whose waterway models are constructed from conservation equations and connected through physical interfaces \cite{OpenHPLUserGuide}. The published Trollheim case study used OpenHPL to model, tune, and validate a real hydropower plant against experimental measurements \cite{Trollheim2019}.

Conceptually,
\[
\boxed{\text{Reservoir.port}\longleftrightarrow\text{Waterway.port}}.
\]
For a fixed reservoir,
\[
H_{\mathrm{port}}=H_0.
\]
For a dynamic reservoir, when port flow is positive into the component,
\[
H_{\mathrm{port}}=H(t),\qquad A(H)\dot H=Q_{\mathrm{in}}+Q_{\mathrm{port}}.
\]
During normal generation at an outlet port, $Q_{\mathrm{port}}<0$.

\section{Implementation in Julia ModelingToolkit / SciML}
ModelingToolkit supports acausal connection semantics. Variables marked with \texttt{connect = Flow} sum to zero over a connection set, while ordinary connector variables are equated \cite{MTKConnect}.

\subsection{Hydraulic Connector}
\begin{lstlisting}
using ModelingToolkit
using ModelingToolkit: t_nounits as t, D_nounits as D

@connector HydraulicPort begin
    H(t), [description = "Hydraulic head [m]"]
    Q(t), [connect = Flow,
           description = "Volumetric flow into component [m^3/s]"]
end
\end{lstlisting}

\subsection{Infinite Reservoir Model}
\begin{lstlisting}
@mtkmodel InfiniteReservoir begin
    @parameters begin
        H0 = 100.0,
        [description = "Constant reservoir head [m]"]
    end
    @components begin
        port = HydraulicPort()
    end
    @equations begin
        port.H ~ H0
    end
end
\end{lstlisting}

\subsection{Constant-Area Dynamic Reservoir}
\begin{lstlisting}
@mtkmodel Reservoir begin
    @parameters begin
        A = 1.0e6,
        [description = "Free-surface area [m^2]"]
        H0 = 100.0,
        [description = "Initial water level [m]"]
    end
    @variables begin
        H(t) = H0,
        [description = "Reservoir water level [m]"]
        Qin(t) = 0.0,
        [description = "External inflow [m^3/s]"]
    end
    @components begin
        port = HydraulicPort()
    end
    @equations begin
        A * D(H) ~ Qin + port.Q
        port.H ~ H
    end
end
\end{lstlisting}

\subsection{Nonlinear Geometry Reservoir}
\begin{lstlisting}
@mtkmodel NonlinearReservoir begin
    @parameters begin
        a0 = 1.0e6
        a1 = 0.0
        a2 = 0.0
        H0 = 100.0
    end
    @variables begin
        H(t) = H0
        Qin(t) = 0.0
        Aeff(t)
    end
    @components begin
        port = HydraulicPort()
    end
    @equations begin
        Aeff ~ a0 + a1*H + a2*H^2
        Aeff * D(H) ~ Qin + port.Q
        port.H ~ H
    end
end
\end{lstlisting}

\section{Model Assembly and Connections}
\begin{lstlisting}
@named reservoir = Reservoir()
@named pipe = RigidPipe()

connections = [
    connect(reservoir.port, pipe.inlet)
]
\end{lstlisting}

The connection contributes
\[
H_{\mathrm{reservoir}}=H_{\mathrm{pipe,in}},
\qquad
Q_{\mathrm{reservoir}}+Q_{\mathrm{pipe,in}}=0.
\]

The future chain can evolve as
\[
\boxed{\texttt{Reservoir}\rightarrow\texttt{RigidPipe}\rightarrow\texttt{SurgeTank}\rightarrow\texttt{Penstock}\rightarrow\texttt{Turbine}}.
\]

\section{Numerical Experiment / Validation}
Consider
\[
A=1000~\si{\meter\squared},\quad
Q_{\mathrm{in}}=10~\si{\meter\cubed\per\second},\quad
Q_{\mathrm{out}}=12~\si{\meter\cubed\per\second}.
\]
Then
\[
\dot H=\frac{10-12}{1000}=-0.002~\si{\meter\per\second}.
\]
After \SI{10}{\second},
\[
H(10)-H(0)=-0.02~\si{\meter},
\]
hence
\[
\boxed{H(10)=H_0-0.02~\si{\meter}}.
\]

\begin{lstlisting}
using ModelingToolkit
using OrdinaryDiffEq

compiled = mtkcompile(model)
prob = ODEProblem(compiled, [], (0.0, 10.0))
sol = solve(prob, Tsit5())
\end{lstlisting}

For larger coupled hydraulic systems, solver selection should be revisited based on stiffness, algebraic structure, and initialization requirements.

\section{Expected Outputs}
For a constant-head reservoir,
\[
H(t)=H_0.
\]
For a dynamic reservoir, useful outputs are
\[
H(t),\quad V(t),\quad Q_{\mathrm{in}}(t),\quad Q_{\mathrm{out}}(t),
\]
and, where useful,
\[
P_h(t)=\rho gQH.
\]

Recommended plots are reservoir level, inflow/outflow, stored volume, and a mass-balance residual. For the explicit outflow convention,
\[
r_m=\dot V-Q_{\mathrm{in}}+Q_{\mathrm{out}},
\qquad
r_m\approx0.
\]
For the connector convention,
\[
\dot V=Q_{\mathrm{in}}+Q_{\mathrm{port}},
\]
and the residual must use the same sign convention.

\section{Engineering Interpretation}
The reservoir establishes a core OpenHPLjl principle:
\[
\boxed{\text{Do not prescribe signal direction unnecessarily.}}
\]
Instead, define physical variables, connector semantics, and conservation laws.

For early SMIB and FCR studies, a useful starting point is
\[
\boxed{H=H_0},
\]
while longer studies can use
\[
\boxed{A(H)\dot H=Q_{\mathrm{in}}-Q_{\mathrm{out}}}.
\]

Thus,
\[
\boxed{\texttt{InfiniteReservoir}\subset\texttt{DynamicReservoir}\subset\texttt{NonlinearReservoir}}
\]
can preserve one hydraulic connection concept while increasing model fidelity.

\section*{Conclusion}
\addcontentsline{toc}{section}{Conclusion}
The reservoir is the natural first component because it introduces
\[
\boxed{\text{physical conservation}+\text{acausal connectors}+\text{DAE formulation}+\text{component reuse}}.
\]
Its fundamental law is
\[
\boxed{\frac{dV}{dt}=Q_{\mathrm{in}}-Q_{\mathrm{out}}},
\]
or
\[
\boxed{A(H)\dot H=Q_{\mathrm{in}}-Q_{\mathrm{out}}}.
\]
The next component is the rigid waterway/rigid pipe, which introduces momentum balance and water inertia.

\begin{thebibliography}{9}
\bibitem{OpenHPLDocs}
OpenHPL documentation, \emph{Open-source hydropower library, Version 3.0.1}.
\url{https://build.openmodelica.org/Documentation/OpenHPL.html}

\bibitem{OpenHPLUserGuide}
OpenHPL, \emph{User's Guide}.
\url{https://build.openmodelica.org/Documentation/OpenHPL.UsersGuide.html}

\bibitem{Trollheim2019}
L. Vytvytskyi and B. Lie,
OpenHPL for Modelling the Trollheim Hydropower Plant,
\emph{Energies}, 12(12), 2303, 2019.
\url{https://doi.org/10.3390/en12122303}

\bibitem{MTKConnect}
SciML, \emph{ModelingToolkit.jl: Acausal Component-Based Modeling}.
\url{https://docs.sciml.ai/ModelingToolkit/stable/tutorials/acausal_components/}
\end{thebibliography}

\end{document}
